% Experiment flow diagram: hierarchical tree showing session structure
% Session -> Segments -> Rounds -> Periods
% Uses TikZ libraries: automata, arrows, positioning, calc (loaded in main.tex)

\begin{figure}[H]
\centering
\resizebox{\textwidth}{!}{%
\begin{tikzpicture}[
    % Node styles
    session/.style={
        rectangle, draw, thick, rounded corners=3pt,
        fill=black!8, minimum width=2.8cm, minimum height=0.9cm,
        font=\bfseries\normalsize
    },
    nochat/.style={
        rectangle, draw, thick, rounded corners=3pt,
        fill=black!12, minimum width=2.6cm, minimum height=0.8cm,
        font=\small, align=center
    },
    chat/.style={
        rectangle, draw, thick, rounded corners=3pt,
        fill=CornflowerBlue!25, minimum width=2.6cm, minimum height=0.8cm,
        font=\small, align=center
    },
    roundnode/.style={
        rectangle, draw, rounded corners=2pt,
        fill=white, minimum width=1.6cm, minimum height=0.7cm,
        font=\small
    },
    periodnode/.style={
        rectangle, draw, rounded corners=2pt,
        fill=white, minimum width=1.4cm, minimum height=0.6cm,
        font=\footnotesize
    },
    ellipsisnode/.style={
        font=\large, inner sep=2pt
    },
    annotation/.style={
        font=\footnotesize\itshape, text=black!70
    },
    arrowed/.style={
        ->, >=stealth, thick
    },
]

% === Session node ===
\node[session] (session) {Part 1};

% === Segment nodes (evenly spaced using absolute x-coordinates) ===
\node[nochat] (seg1) at (-6, -2.2) {Segment 1\\(No Chat)};
\node[nochat] (seg2) at (-2, -2.2) {Segment 2\\(No Chat)};
\node[chat]   (seg3) at ( 2, -2.2) {Segment 3\\(Chat)};
\node[chat]   (seg4) at ( 6, -2.2) {Segment 4\\(Chat)};

% === Session -> Segment edges ===
\draw[arrowed] (session) -- (seg1);
\draw[arrowed] (session) -- (seg2);
\draw[arrowed] (session) -- (seg3);
\draw[arrowed] (session) -- (seg4);

% === Vertical dots under segments without expanded subtree ===
\node[annotation] at (-6, -3.5) {$\vdots$};
\node[annotation] at ( 2, -3.5) {$\vdots$};
\node[annotation] at ( 6, -3.5) {$\vdots$};

% === Round nodes (under Segment 2) ===
\node[roundnode] (r1) at (-4.2, -4.6) {Round 1};
\node[roundnode] (r2) at (-2.0, -4.6) {Round 2};
\node[ellipsisnode] (rdots) at (-0.3, -4.6) {$\cdots$};
\node[roundnode] (rR) at (1.2, -4.6) {Round $R$};

% === Segment -> Round edges ===
\draw[arrowed] (seg2) -- (r1);
\draw[arrowed] (seg2) -- (r2);
\draw[arrowed] (seg2) -- (rdots);
\draw[arrowed] (seg2) -- (rR);

% === Geometric annotation for rounds ===
\node[annotation, right=0.4cm of rR]
    {$R \sim \text{Geometric}(0.125)$};

% === Period nodes (under Round 1) ===
\node[periodnode] (p1) at (-6.2, -7.0) {Period 1};
\node[periodnode] (p2) at (-4.4, -7.0) {Period 2};
\node[ellipsisnode] (pdots) at (-3.0, -7.0) {$\cdots$};
\node[periodnode] (pT) at (-1.6, -7.0) {Period $T$};

% === Round -> Period edges ===
\draw[arrowed] (r1) -- (p1);
\draw[arrowed] (r1) -- (p2);
\draw[arrowed] (r1) -- (pdots);
\draw[arrowed] (r1) -- (pT);

% === Geometric annotation for periods ===
\node[annotation, right=0.4cm of pT]
    {$T \sim \text{Geometric}(0.125)$};

\end{tikzpicture}%
}% end resizebox
\caption{Hierarchical structure of the experiment.}
\label{fig:experiment_flow}
\end{figure}
